\begin{thebibliography}{99}
\raggedright
\bibitem{TP3} G. Yang. \emph{Tensor Programs III: Neural Matrix Laws}.
\href{https://arxiv.org/abs/2009.10685}{arXiv:2009.10685}, 2020.
Fixed-computation Gaussian response theorem: Theorem 2.10, Box 1,
Remarks 2.11--2.12.
\bibitem{TP4} G. Yang and E. J. Hu. \emph{Tensor Programs IV: Feature Learning in Infinite-Width Neural Networks}.
\href{https://arxiv.org/abs/2011.14522}{arXiv:2011.14522}, 2020.
\bibitem{TP4b} G. Yang and E. Littwin. \emph{Tensor Programs IVb: Adaptive Optimization in the Infinite-Width Limit}.
\href{https://arxiv.org/abs/2308.01814}{arXiv:2308.01814}, 2023.
Operator semantics: Definition 2.6.7 and Propositions 2.6.8--2.6.9.
\bibitem{TP5} G. Yang et al. \emph{Tensor Programs V: Tuning Large Neural Networks via Zero-Shot Hyperparameter Transfer}.
\href{https://arxiv.org/abs/2203.03466}{arXiv:2203.03466}, 2022.
\bibitem{ProjectBorel} PDE project.
\emph{Autonomous Deep Feature Learning: Theorem and Program},
21 August 2026, Section 10 added 23 August 2026.
Source of the one-input, initialization-jet Borel reconstruction conjecture
and its compact-time variant. Project source file:
\nolinkurl{AUTONOMOUS_DEEP_FEATURE_LEARNING_THEOREM_AND_PROGRAM_2026-08-21.md}.
\bibitem{ProjectMonograph} PDE project.
\emph{Finite Causal Neural PDE Master Monograph}, version 2.2,
31 July 2026, Section 5.7.6.
Discussion of flat initialization jets and reconstruction.
Project source file:
\nolinkurl{FINITE_CAUSAL_NEURAL_PDE_MASTER_MONOGRAPH_v2.2_2026-07-31.md}.
\bibitem{LastraMalekSanz} A. Lastra, S. Malek and J. Sanz.
\emph{Summability in general Carleman ultraholomorphic classes}.
\href{https://arxiv.org/abs/1402.1669}{arXiv:1402.1669}, 2014.
\bibitem{AllenChui} G. D. Allen, C. K. Chui, W. R. Madych,
F. J. Narcowich and P. W. Smith.
\emph{Pad\'e approximation and Gaussian quadrature}.
Bulletin of the Australian Mathematical Society 11 (1974), 63--69.
\href{https://doi.org/10.1017/S0004972700043641}{doi:10.1017/S0004972700043641}.
\bibitem{SimonMoments} B. Simon.
\emph{The Classical Moment Problem as a Self-Adjoint Finite Difference Operator}.
Advances in Mathematics 137 (1998), 82--203.
\href{https://arxiv.org/abs/math-ph/9906008}{arXiv:math-ph/9906008}.
\bibitem{BorghiWeniger} R. Borghi and E. J. Weniger.
\emph{Convergence Analysis of the Summation of the Euler Series by Pad\'e
Approximants and the Delta Transformation}.
\href{https://arxiv.org/abs/1405.2474}{arXiv:1405.2474}, 2014.
\bibitem{LiuPego} J.-G. Liu and R. L. Pego.
\emph{On generating functions of Hausdorff moment sequences}.
\href{https://arxiv.org/abs/1401.8052}{arXiv:1401.8052}, 2014.
\bibitem{Cacoq} J. Cacoq, B. de la Calle Ysern and G. L\'opez Lagomasino.
\emph{Incomplete Pad\'e approximation and convergence of row sequences of
Hermite--Pad\'e approximants}.
Journal of Approximation Theory 170 (2013), 59--77.
\href{https://arxiv.org/abs/1111.2774}{arXiv:1111.2774}.
\bibitem{Stahl} H. Stahl.
\emph{The Convergence of Pad\'e Approximants to Functions with Branch Points}.
Journal of Approximation Theory 91 (1997), 139--204.
\href{https://doi.org/10.1006/jath.1997.3141}{doi:10.1006/jath.1997.3141}.
\bibitem{Lubinsky} D. S. Lubinsky.
\emph{Rogers--Ramanujan and the Baker--Gammel--Wills (Pad\'e) conjecture}.
Annals of Mathematics 157 (2003), 847--889.
\href{https://annals.math.princeton.edu/2003/157-3/p04}{Journal article}.
\bibitem{Buslaev} V. I. Buslaev.
\emph{On the Baker--Gammel--Wills conjecture in the theory of Pad\'e approximants}.
Sbornik: Mathematics 193:6 (2002), 811--823.
\href{https://www.mathnet.ru/eng/sm658}{Journal article}.
\bibitem{BeckermannMatos} B. Beckermann and A. C. Matos.
\emph{Algebraic properties of robust Pad\'e approximants}.
\href{https://arxiv.org/abs/1310.2438}{arXiv:1310.2438}, 2013.
\bibitem{CostinDunne} O. Costin and G. V. Dunne.
\emph{Uniformization and Constructive Analytic Continuation of Taylor Series}.
\href{https://arxiv.org/abs/2009.01962}{arXiv:2009.01962}, 2020.
\end{thebibliography}
